% BHU PYQ -- Transcribed & Corrected 100% Accurate (English Version)
% Paper: SOCMD11: Invitation to Sociology
% Exam: Under Graduate Multidisciplinary Course, Semester I Examination 2024-25 (NEP)
% Category: Multidisciplinary Course (MDC)

\documentclass[11pt,a4paper]{article}
\usepackage[a4paper,top=2.5cm,bottom=2.5cm,left=2.5cm,right=2.5cm,headheight=14pt]{geometry}
\usepackage{amsmath,amssymb}
\usepackage[T1]{fontenc}
\usepackage[utf8]{inputenc}
\usepackage{lmodern,microtype,fancyhdr,enumitem,hyperref,lastpage,parskip}

\hypersetup{colorlinks=true,urlcolor=black,linkcolor=black,pdftitle={SOCMD11: Invitation to Sociology -- BHU Sem I 2024-25 (NEP MDC)}}
\pagestyle{fancy}\fancyhf{}
\fancyhead[L]{\small\textit{Banaras Hindu University}}
\fancyhead[R]{\small\textit{SOCMD11: Invitation to Sociology (MDC)}}
\fancyfoot[C]{\small Page \thepage\ of \pageref{LastPage}}

\newlist{parts}{enumerate}{2}
\setlist[parts,1]{label=(\roman*),leftmargin=*,itemsep=4pt,topsep=2pt}
\newcommand{\pts}[1]{\hfill\textit{[#1]}}

\begin{document}

\begin{center}
  {\large\bfseries BANARAS HINDU UNIVERSITY}\\[2pt]
  {\normalsize Under Graduate Multidisciplinary Course, Semester I Examination 2024-25 (NEP)}\\[6pt]
  \rule{0.8\linewidth}{0.5pt}\\[6pt]
  {\large\bfseries DEPARTMENT OF SOCIOLOGY}\\[4pt]
  {\large\bfseries Multidisciplinary Course (MDC)}\\[4pt]
  {\large\bfseries Paper Code: SOCMD11 --- Invitation to Sociology}\\[4pt]
  {\normalsize Time: 2:30 Hours\hfill Maximum Marks: 70}
\end{center}

\rule{\linewidth}{0.4pt}\smallskip
\noindent\textit{Instruction: Write your Roll No. at the top immediately on the receipt of this question paper.}\\[2pt]
\noindent\textit{Note: The question paper is divided into three sections: Section-A (Objective/Short type, 2 marks each), Section-B (Short answer type, 10 marks each, attempt 3), Section-C (Long answer type, 15 marks each, attempt 2).}
\bigskip

\begin{center}
  \textbf{\large SECTION --- A}
\end{center}
\textit{(Objective / Short Answer Questions --- 2 Marks each, max 50 words)}
\bigskip

\noindent\textbf{Question 1} \textit{Answer the following:}\pts{$5 \times 2 = 10$}
\begin{parts}
  \item What do you mean by `Subjectivity' and `Objectivity'?
  \item What do you mean by `Cultural Diffusion'?
  \item Define `Social Structure'.
  \item What do you mean by `Social Transformation'?
  \item What do you mean by `Cultural Lag'?
\end{parts}

\bigskip
\begin{center}
  \textbf{\large SECTION --- B}
\end{center}
\textit{(Short Answer Type Questions --- 10 Marks each, max 250 words. Attempt any THREE questions)}
\bigskip

\noindent\textbf{Question 2}\pts{10} \\
Explain the relationship between sociology and common sense.

\begin{center}\textbf{OR}\end{center}

What are the major criticisms of common sense as a tool for understanding social phenomena?

\medskip\rule{\linewidth}{0.2pt}\medskip

\noindent\textbf{Question 3}\pts{10} \\
Explain the concept of Sociological Imagination and its relevance in understanding society. Use examples to illustrate your answer.

\begin{center}\textbf{OR}\end{center}

Discuss the relationship between culture and social structure. How do culture and social institutions influence each other?

\medskip\rule{\linewidth}{0.2pt}\medskip

\noindent\textbf{Question 4}\pts{10} \\
Discuss the concept of `role performance' and its significance in social interaction.

\begin{center}\textbf{OR}\end{center}

What is the role of culture in shaping social interaction?

\bigskip
\begin{center}
  \textbf{\large SECTION --- C}
\end{center}
\textit{(Long Answer Type Questions --- 15 Marks each. Attempt any TWO questions)}
\bigskip

\noindent\textbf{Question 5}\pts{15} \\
What do you mean by Institution? Discuss its evolution.

\begin{center}\textbf{OR}\end{center}

How do individuals construct their identities in relation to societal expectations and cultural norms?

\medskip\rule{\linewidth}{0.2pt}\medskip

\noindent\textbf{Question 6}\pts{15} \\
Sociologists often argue that everyday life is both ordinary and deeply structured. Discuss this statement using examples from both urban and rural settings.

\begin{center}\textbf{OR}\end{center}

Define Social Change. Discuss the various factors responsible for social change.

\vfill\rule{\linewidth}{0.4pt}
\begin{center}\small
\href{https://bhu-exam.vercel.app/}{Click here to visit our website}\\[4pt]
\href{https://me-aryan.vercel.app/}{Click here to visit our contributor}
\end{center}

\end{document}
